\section{Conjunto de datos}
    \subsection{Origen y licencia}
        El conjunto \textit{Resultados únicos Saber 11} es publicado por el
        Instituto Colombiano para la Evaluación de la Educación (ICFES) en el
        Portal Nacional de Datos Abiertos de Colombia, bajo el identificador
        \texttt{kgxf-xxbe} y con licencia Creative Commons Attribution-ShareAlike
        4.0 (CC BY-SA 4.0), que permite su uso, redistribución y adaptación con
        atribución \cite{icfes2023resultados}. Reúne 7,1 millones de registros
        anónimos de las aplicaciones realizadas entre 2010 y 2022, uno por
        estudiante y periodo, con 51 columnas, y su última actualización data de
        agosto de 2023. El portal expone los datos mediante la API de Socrata, lo
        que permite descargar un periodo específico con un filtro en la URL sin
        necesidad de registro.

    \subsection{Periodo seleccionado y tamaño}
        Se trabaja con el periodo 2022-4, correspondiente a la aplicación del
        calendario A del segundo semestre de 2022, que es la que presenta la gran
        mayoría de colegios del país y el periodo más reciente y numeroso del
        conjunto. La descarga pesa 560 MB y contiene 1.065.888 filas; como se
        muestra en la Sección~V, la mitad de ellas son duplicados exactos y el
        conjunto de trabajo queda en 532.792 estudiantes. La Tabla~\ref{dataset_summary}
        resume sus características.

        \begin{table}[H]
            \centering
            \caption{Resumen del conjunto de datos}
            \label{dataset_summary}
            \begin{tabular}{ll}
                \toprule
                \textbf{Aspecto} & \textbf{Valor} \\
                \midrule
                Fuente & datos.gov.co, conjunto \texttt{kgxf-xxbe} \\
                Publicador & ICFES \\
                Licencia & CC BY-SA 4.0 \\
                Periodo & 2022-4 (calendario A, segundo semestre) \\
                Filas descargadas & 1.065.888 \\
                Filas únicas & 532.792 \\
                Columnas & 51 \\
                Variable objetivo & \texttt{punt\_global} (0 a 500) \\
                Tamaño en disco & 560 MB (CSV) \\
                Tamaño en memoria & 131 MB (tipos categóricos) \\
                \bottomrule
            \end{tabular}
        \end{table}

    \subsection{Variables}
        Las 51 columnas se agrupan en cinco bloques según su prefijo, descritos
        en la Tabla~\ref{variables}. Todas llegan como texto; en el cuaderno las
        categóricas se leen como \texttt{category}, los puntajes como numéricos y
        los identificadores como cadena.

        \begin{table}[H]
            \centering
            \caption{Bloques de variables del conjunto}
            \label{variables}
            \begin{tabular}{p{2.1cm}p{5.6cm}}
                \toprule
                \textbf{Bloque} & \textbf{Contenido} \\
                \midrule
                Identificación (5) & Periodo, identificador anónimo
                (\texttt{estu\_consecutivo}), tipo de documento, condición del
                evaluado y estado del resultado. \\
                Colegio, \texttt{cole\_*} (17) & Códigos DANE e ICFES, nombres
                del establecimiento y la sede, naturaleza (oficial o no oficial),
                área (urbana o rural), calendario, carácter (académico o técnico),
                jornada, género de la población, bilingüismo y ubicación
                (departamento y municipio). \\
                Estudiante, \texttt{estu\_*} (13) & Género, fecha de nacimiento,
                nacionalidad, lugar de residencia, lugar de presentación de la
                prueba y privación de la libertad. \\
                Hogar, \texttt{fami\_*} (9) & Estrato de la vivienda, máximo
                nivel educativo de la madre y del padre, personas y cuartos del
                hogar y tenencia de internet, computador, automóvil y lavadora;
                provienen del cuestionario que el estudiante diligencia al
                inscribirse. \\
                Resultados (7) & Puntaje en lectura crítica, matemáticas,
                sociales y ciudadanas, ciencias naturales e inglés (0 a 100),
                nivel de desempeño en inglés y puntaje global (0 a 500). \\
                \bottomrule
            \end{tabular}
        \end{table}

        La variable objetivo es el puntaje global. Las variables explicativas
        candidatas son las de los bloques de estudiante, hogar y colegio, más el
        departamento y el municipio del colegio; los resultados por área se
        excluyen por la razón que se verifica en la Sección~V.
